\magnification=\magstep1
\input notesmac.tex
\nopagenumbers

\centerline{\titlefont Tech and Tools, Tutorial 2}
\bigskip
\bigskip
\line{Dr.~WANG Jian-Sheng \hfill 27 Sep 1999}
\bigskip
\bigskip
{\sl I'll divided into $n$ groups, for $n=1, 2, 3, \cdots$.}
\bigskip

%----------------------------------------------------------------------- 

\problem  After specifying the following assignments and rules,
\ttbg
f[2] = 3
f[u] := 2u^2
f[v_] := 1/v
u = 2
\tted
what is {\tt f[u]}.


\bigskip
\problem  Explain what the following function {\tt func[]} can 
accomplish. Give an example to illustrate the use of {\tt func[]}.
\ttbg
func[l_] := ff[l, {}];
ff[ {}, r_] := r;
ff[ {h_, t___}, {r___} ] := ff[ {t}, {r, h} ] /; h>0;
ff[ {h_, t___}, {r___} ] := ff[{t}, {r}] /; (h<=0||Head[h]==Symbol)
\tted

\bigskip
\problem What does this program do?
\ttbg
runEncode[lst_List] :=
  Map[ {#,1}&, lst] //. 
     {u___, {v_,r_}, {v_, s_}, w___} -> {u, {v, r+s}, w}
\tted


\bigskip
\problem How does one read a data file representing a 4 by 4 matrix
\ttbg
  0.3 -3.0 0.3  1.023
  0.1  0.2 1.1  0.003
  5.0  2.3 2.3  0.1
  0.3  0.1 4.4  0.005
\tted
into Mathematica as a (Mathematica) matrix (list of list) for further
manipulation (such as solving a linear equation)?

\item{} Conversely, how does one write a set of data suitable for C or Fortran
to read? 

\bye
